% 图 74.11 —— 由 `checks/emit_hand.py` 自动拼出（probe 精测 + prims2 图元分解）
%%
% ⚠ 这是**稿子不是成品**：跑 `checks/checkhand.py 74.11` 看指标 + 看叠合图。
%%
%% ⚠ 待核：
%   · 本图 probe 报出阴影族 1 个 —— **本脚本不生成阴影**（要照原书手写）；prims2 的 strokes 里可能含阴影线，核对时留意别画重
%   · 可疑字形 2 项（空心圈 ⊙ 常被认成字母 o）—— 见文件末
%%
\documentclass[12pt]{standalone}
\usepackage{fontspec}
\setsansfont{Arial}
\newfontfamily\mfallback{DejaVu Sans}
\usepackage{tikz}
\begin{document}
\begin{tikzpicture}[x=1pt, y=-1pt, line width=6.0pt, line cap=round, line join=round]

  %% ════ 填充区（prims2 的 fills）════

  %% ════ 直线段（prims2 的 strokes）════
  \draw[line width=6.0pt] (1052.0,200.0) -- (1051.0,192.0);
  \draw[line width=4.5pt] (1051.0,201.0) -- (1039.0,190.0);
  \draw[line width=5.0pt] (843.0,278.0) -- (840.0,307.0);
  \draw[line width=4.0pt] (482.0,195.0) -- (481.0,167.0);
  \draw[line width=5.0pt] (638.0,166.0) -- (482.0,166.0);
  \draw[line width=5.0pt] (1066.0,207.0) -- (1071.4,205.6);
  \draw[line width=5.0pt] (850.6,206.4) -- (832.0,223.0);
  \draw[line width=4.9pt] (1065.0,208.0) -- (1063.6,212.4);
  \draw[line width=5.0pt] (1063.6,212.4) -- (966.0,273.0);
  \draw[line width=5.0pt] (856.0,275.0) -- (845.0,276.0);
  \draw[line width=5.0pt] (975.0,354.0) -- (986.0,356.0);
  \draw[line width=7.0pt] (1080.0,109.0) -- (1075.0,105.0);
  \draw[line width=4.0pt] (905.1,191.0) -- (889.6,204.4);
  \draw[line width=4.5pt] (889.6,204.4) -- (867.0,229.0);
  \draw[line width=5.0pt] (49.0,190.0) -- (49.8,299.8);
  \draw[line width=5.0pt] (49.8,299.8) -- (185.0,300.0);
  \draw[line width=6.0pt] (475.0,191.0) -- (480.0,197.0);
  \draw[line width=4.0pt] (800.0,170.0) -- (817.0,170.0);
  \draw[line width=6.1pt] (639.0,165.0) -- (636.0,160.0);
  \draw[line width=4.0pt] (883.0,262.0) -- (870.0,271.0);
  \draw[line width=5.0pt] (872.0,305.0) -- (886.0,296.0);
  \draw[line width=9.0pt] (1090.0,108.0) -- (1082.0,111.0);
  \draw[line width=5.0pt] (640.0,166.0) -- (724.0,165.0);
  \draw[line width=6.0pt] (184.0,50.0) -- (187.0,45.0);
  \draw[line width=5.0pt] (1018.4,79.6) -- (1030.2,74.0);
  \draw[line width=5.0pt] (1030.2,74.0) -- (1038.8,75.0);
  \draw[line width=5.0pt] (469.0,139.0) -- (704.0,137.0);
  \draw[line width=5.0pt] (704.0,137.0) -- (710.8,140.0);
  \draw[line width=6.1pt] (185.0,38.0) -- (188.0,43.0);
  \draw[line width=5.0pt] (908.0,247.0) -- (893.0,257.0);
  \draw[line width=5.0pt] (712.5,228.5) -- (701.7,234.0);
  \draw[line width=5.0pt] (701.7,234.0) -- (467.0,233.0);
  \draw[line width=6.1pt] (1045.0,203.0) -- (1050.0,202.0);
  \draw[line width=6.0pt] (823.0,262.0) -- (830.0,268.0);
  \draw[line width=5.0pt] (727.0,190.0) -- (724.0,165.0);
  \draw[line width=5.2pt] (840.0,308.0) -- (843.0,312.0);
  \draw[line width=7.0pt] (1053.0,202.0) -- (1065.0,207.0);
  \draw[line width=4.0pt] (186.0,44.0) -- (50.7,44.3);
  \draw[line width=5.0pt] (50.7,44.3) -- (49.0,187.0);
  \draw[line width=4.0pt] (911.0,280.0) -- (897.0,290.0);
  \draw[line width=4.0pt] (1038.0,188.0) -- (1031.0,172.0);
  \draw[line width=6.0pt] (468.0,139.0) -- (449.7,141.3);
  \draw[line width=5.0pt] (449.1,232.0) -- (466.0,233.0);
  \draw[line width=5.0pt] (981.0,363.0) -- (986.0,357.0);
  \draw[line width=4.0pt] (799.0,164.0) -- (819.0,163.0);
  \draw[line width=4.0pt] (300.0,150.0) -- (299.7,298.3);
  \draw[line width=5.0pt] (299.7,298.3) -- (188.0,301.0);
  \draw[line width=5.0pt] (844.0,312.0) -- (859.0,309.0);
  \draw[line width=4.0pt] (934.0,264.0) -- (920.0,272.0);
  \draw[line width=5.0pt] (481.0,166.0) -- (469.0,140.0);
  \draw[line width=4.0pt] (928.0,233.0) -- (919.0,240.0);
  \draw[line width=6.0pt] (42.0,185.0) -- (48.0,189.0);
  \draw[line width=6.0pt] (50.0,188.0) -- (55.0,185.0);
  \draw[line width=4.0pt] (962.0,284.0) -- (960.0,289.0);
  \draw[line width=4.0pt] (300.0,147.0) -- (301.0,47.2);
  \draw[line width=3.0pt] (301.0,47.2) -- (299.0,45.0);
  \draw[line width=5.0pt] (299.0,45.0) -- (189.0,44.0);
  \draw[line width=6.1pt] (301.0,148.0) -- (306.0,151.0);

  %% ════ 曲线（prims2 的 curves —— 三段式，坐标同为 (y,x)）════
  \draw[line width=5.0pt] (1071.4,205.6) .. controls (1109.6,170.0) and (1125.3,111.8) .. (1104.9,63.9);
  \draw[line width=5.0pt] (1104.9,63.9) .. controls (1084.4,25.6) and (1034.2,14.5) .. (995.2,27.0);
  \draw[line width=5.0pt] (995.2,27.0) .. controls (950.6,45.1) and (937.6,94.6) .. (925.0,134.6);
  \draw[line width=5.0pt] (925.0,134.6) .. controls (908.4,164.8) and (879.8,189.7) .. (850.6,206.4);
  \draw[line width=5.0pt] (832.0,223.0) .. controls (825.5,234.0) and (818.1,247.3) .. (822.0,261.0);
  \draw[line width=5.0pt] (1080.0,109.0) .. controls (1085.2,94.2) and (1079.3,76.8) .. (1071.9,63.9);
  \draw[line width=5.0pt] (1071.9,63.9) .. controls (1047.9,35.4) and (1001.0,40.6) .. (977.2,65.8);
  \draw[line width=4.0pt] (977.2,65.8) .. controls (949.3,105.8) and (951.4,161.8) .. (905.1,191.0);
  \draw[line width=5.0pt] (867.0,229.0) .. controls (858.5,243.9) and (847.1,258.5) .. (845.0,276.0);
  \draw[line width=4.0pt] (1039.0,188.0) .. controls (1061.4,167.0) and (1078.4,141.7) .. (1081.0,112.0);
  \draw[line width=5.0pt] (989.0,191.0) .. controls (988.5,152.0) and (990.6,109.0) .. (1018.4,79.6);
  \draw[line width=4.0pt] (1038.8,75.0) .. controls (1075.9,98.4) and (1049.5,144.7) .. (1031.0,171.0);
  \draw[line width=5.0pt] (710.8,140.0) .. controls (718.4,145.3) and (725.2,156.0) .. (724.0,165.0);
  \draw[line width=5.0pt] (726.0,193.0) .. controls (725.3,206.0) and (722.5,219.2) .. (712.5,228.5);
  \draw[line width=5.0pt] (1030.0,172.0) .. controls (1002.8,198.8) and (968.4,211.2) .. (937.0,230.0);
  \draw[line width=5.0pt] (839.0,307.0) .. controls (826.0,297.8) and (819.6,278.4) .. (821.0,263.0);
  \draw[line width=5.0pt] (449.7,141.3) .. controls (428.8,165.1) and (427.7,208.3) .. (449.1,232.0);
  \draw[line width=5.0pt] (950.0,249.0) .. controls (981.6,233.5) and (1011.0,214.9) .. (1037.0,190.0);
  \draw[line width=5.0pt] (843.0,313.0) .. controls (871.5,357.5) and (931.6,391.0) .. (981.0,363.0);
  \draw[line width=5.0pt] (987.0,356.0) .. controls (1017.2,335.9) and (1019.8,292.5) .. (1013.0,260.0);
  \draw[line width=5.0pt] (480.0,197.0) .. controls (479.4,210.0) and (474.0,221.5) .. (467.0,232.0);

  %% ════ 实心点 ════
  \fill (485.5,199.0) circle (5.6);
  \fill (187.5,306.0) circle (4.6);

  %% ════ 标签（probe 直接给的 \node 行）════
  %% ⚠ 全部补了 fill=white：文字在最上层，否则与曲线交叉干扰
     \node[anchor=center,inner sep=0pt,fill=white,font=\fontsize{32.0}{40.0}\selectfont] at (155.1,27.1) {\textsf{\textbf{\textit{a}}}};   % box(146,17,16,17)
     \node[anchor=center,inner sep=0pt,fill=white,font=\fontsize{32.0}{40.0}\selectfont] at (1097.1,132.1) {\textsf{\textbf{\textit{a}}}};   % box(1088,122,16,17)
     \node[anchor=center,inner sep=0pt,fill=white,font=\fontsize{30.7}{38.3}\selectfont] at (30.8,169.8) {\textsf{\textbf{\textit{b}}}};   % box(21,157,17,23)
     \node[anchor=center,inner sep=0pt,fill=white,font=\fontsize{30.0}{37.5}\selectfont] at (327.8,170.2) {\textsf{\textbf{\textit{b}}}};   % box(318,158,17,22)
     \node[anchor=center,inner sep=0pt,fill=white,font=\fontsize{41.3}{51.6}\selectfont] at (388.5,189.2) {{\mfallback\textbf{−}}};   % box(368,180,21,6)
     \node[anchor=center,inner sep=0pt,fill=white,font=\fontsize{30.0}{37.6}\selectfont] at (595.5,189.5) {\textsf{\textbf{\textit{a}}}};   % box(587,180,15,16)
     \node[anchor=center,inner sep=0pt,fill=white,font=\fontsize{41.3}{51.6}\selectfont] at (388.5,195.2) {{\mfallback\textbf{−}}};   % box(368,186,21,6)
     \node[anchor=center,inner sep=0pt,fill=white,font=\fontsize{30.0}{37.5}\selectfont] at (749.8,198.2) {\textsf{\textbf{\textit{b}}}};   % box(740,186,17,22)
     \node[anchor=center,inner sep=0pt,fill=white,font=\fontsize{30.0}{37.5}\selectfont] at (502.8,206.2) {\textsf{\textbf{\textit{b}}}};   % box(493,194,17,22)
     \node[anchor=center,inner sep=0pt,fill=white,font=\fontsize{30.0}{37.5}\selectfont] at (1100.8,209.2) {\textsf{\textbf{\textit{b}}}};   % box(1091,197,17,22)
     \node[anchor=center,inner sep=0pt,fill=white,font=\fontsize{22.5}{28.1}\selectfont] at (956.4,310.4) {\textsf{\textbf{(}}};   % box(953,300,6,20)
     \node[anchor=center,inner sep=0pt,fill=white,font=\fontsize{32.9}{41.1}\selectfont] at (152.1,326.6) {\textsf{\textbf{\textit{a}}}};   % box(143,316,16,18)
     \node[anchor=center,inner sep=0pt,fill=white,font=\fontsize{61.4}{76.7}\selectfont] at (959.5,354.3) {\textsf{\textbf{′}}};   % box(955,330,11,18)
     \node[anchor=center,inner sep=0pt,fill=white,font=\fontsize{30.7}{38.4}\selectfont] at (1040.5,375.5) {\textsf{\textbf{\textit{K}}}};   % box(1028,363,22,23)

  % 钉死画布：否则超出的图元/控制点会把页面撑大，1:1 回比就失准
  \pgfresetboundingbox
  \path[use as bounding box] (0,0) rectangle (1130,396);
\end{tikzpicture}
\end{document}

% ⚠ 待核清单（probe 拿不准，人眼过一遍）：
%   · 未识别/跳过：⑤ 跳过/未识别的字形
%   · 未识别/跳过：box(798,168,22,6)
